\documentclass[../main.tex]{subfiles}

\begin{document}

\chapter{Induksi Matematika dan Penerapannya pada Algoritma Rekursif}

\begin{subcpmk}
  \item Mahasiswa mampu mengaplikasikan prinsip penalaran deduktif iteratif Induksi Matematika untuk mengevaluasi parameter rumus dan membuktikan universal logikal fungsi kebenaran deret atau komputasi \textit{looping array}/\textit{recursive}.
\end{subcpmk}

\section{Prinsip Dasar Induksi Matematis di Ranah Komputasional Logik}

Banyak sekali spesifikasi perancangan perangkat lunak \textit{(Program specifications)} yang dibangun dari fungsi atau argumen algoritmik berspesifikasi pengulangan/komputasi berjalan numerik kontinu berseri (\textit{Loops} dan Deret algoritmis dari indeks ke-0 menyusur ke $n$ iterasi). Bilamana parameter fungsi algoritmis kita mengklaim sifat operasional var fungsional $P(n)$ bernilai universal *VALID* bagi SETIAP elemen domain data himpunan komputasi Bilangan Asli $n \geq 1$, menguji setiap himpunan datanya via metode enumeratif \textit{brute-force} tidaklah mungkin diselesaikan (karena himpunannya \textit{Infinite/tak hingga} batas jumlah indeks var numeriknya !).

Untuk mensertifikasi validitas fungsional logis tak terhingga iteratif var array struktur komputasinya, \textit{Software Architect} mempergunakan senjata penalaran rasional komputatif mutlak murni: \textbf{INDUKSI MATEMATIKA (\textit{Mathematical Induction})}.
Mirip fenomena Efek Domino mekanik; Anda tak perlu menjatuhkan sendiri sekuens domino satu per satu secara manual. \textit{Cukup garansi ke kebenaran absolut validitas bahwa ubin domino nomor $1$ berhasil dirobohkan komputasional sistem logik komparator (*Base Case*)}, dan validasi \textit{Hukum Fisika fungsional bahwa JIKA sebuah domino fungsional iteratif berseri logikal rasional roboh/valid (*asumsi ke$-k$*), MAKA ia MESTI (dipastikan rasional logis formula deduktif memvalidasi secara mutlak absolut memvalidasi!) menabrak dan memicu ubin logis beriteratif indeks bersebelahannya ke-$(k+1)$ ikut roboh sukses fungsional (*Inductive Step/Transition Rule*)}. Seketika seluruh jajaran infinit sekuens domino array \textit{loop/rekursif}-nya bakal runtuh fungsional berhasil mutlak terbukti T/True rasional logis!

\section{Dua Langkah Mekanik Algoritmis Induksi Formal}

\subsection{1. Langkah \textit{Base Case} (Basis Induksi Pembuktian)}
Kita inisialisasi iterasi komputasinya. Buktikan secara manual/substitusi *direct* nilai rumus/fungsional kalkulasinya bahwa \textbf{$P(1)$} (Atau parameter awal *index var loop* algoritma terkecil var \textit{array} komputasinya, misal 0/1) adalah bernilai mutlak BENAR \textit{/ True Value logic}.

\subsection{2. Langkah \textit{Inductive Step} (Hipotesis dan Langkah Induksi)}
Di sesi sirkulasi fungsional komputasi algoritmik rasional perumusannya:
\begin{itemize}
    \item \textbf{Asumsi Induktif Mutlak:} (Tanpa perlu buktikan ulang) Asumsikan bahwa iterasi $P(k)$ berlaku *Valid komputatif True* untuk seluruh rentang logikal $k \geq 1$ sebarang komputatif. (Asumsikan argumen Domino pada iterasi angka array index $k$ roboh mutlak sukses).
    \item \textbf{Eksekusi Langkah Pembuktian \textit{Target}:} Atas spesifikasi modal kelebihan fungsi matematis \textit{asumsi $P(k)$} tadi, rakitlah \textit{logical string}/\textit{algebra manipulation} formula parameter rumus, yang menunjukkan konklusif bahwa argumen predikat variabel logika komputasi \textit{array} $P(k+1)$ OTOMATIS IKUTAN BER-STATUS BENAR (\textit{VALID FUNGSIONAL ALGORITMIK T/TRUE}).
\end{itemize}
Bilamana kedua gerbang \textit{test check komparator array logis} induktif di fungsional \textit{Base Case} dan fungsi \textit{Step} berhasil disimulasikan sah rasional tanpa error kontradiksi kalkulasi, *Selesailah kompilasi pembuktian*. Klaim fungsional Anda disahkan mutlak kebenarannya terhadap SELURUH ELEMEN INFINITE INTEGER. \textit{Q.E.D}.

\begin{contoh}
\textbf{Pembuktian Deret Integer Loop Algoritma}
Kompilator mengklaim fungsi Penjumlahan *Array loop* \textit{Sum} $(1 + 2 + 3 \dots + n)$ komputabilitas rasional evaluatif formulanya diganti/dioptimisasi dengan struktur fungsi variabel rumus mutlak $= \frac{n(n+1)}{2}$. Buktikan!
\\
\textbf{Basis Langkah Komputasi:} $(n=1)$
Sebelah Kiri: $= 1$. \textit{Logic} Sebelah Kanan Terminal: $\frac{1(1+1)}{2} = 1$. \textit{Evaluatif Persamaan \textbf{True/Sama}}. (Basis Terverifikasi Logikal Sah).
\\
\textbf{Langkah Induksi Loop Algoritma:}
- Diasumsikan untuk tipe var ke-$k$, maka $(1 + 2 \dots + k) = \frac{k(k+1)}{2}$ [Statusnya diasumsikan *Sah Absolut/T*].
- Akan disintesis dan dipersyaratkan dibuktikan terminal fungsional berikutnya di var ke-$(k+1)$:  $(1+2\dots+k) + (k+1) $ wajib \textit{ekivalen logikal sama terminal resultan formulasinya formula manipulasif komparatornya dengan substitusinya} $= \frac{(k+1)((k+1)+1)}{2} = \frac{(k+1)(k+2)}{2}$.
- \textit{Sintesis Pembuktian Modifikasi fungsional dari kiri}:
  Bagian iterasi pertama $(1+2\dots+k)$ kita ganti mutlak var logikalnya sesuai Hipotesis komputasinya $\to$ $\frac{k(k+1)}{2} + (k+1)$.
- Kita perhalus spesifikasi komputasinya secara analitik: $= \frac{k(k+1)}{2} + \frac{2(k+1)}{2}$.
- Gabungkan parameternya $ = \frac{(k+1)(k+2)}{2} $.
- Coba kita verifikasi hasil \textit{Kiri terminal final}= \textbf{TERBUKTI PERSIS SAMA IDENTIK LOGIKAL} dan rasional komputasinya sinkron berpadu mutlak akurasinya dengan desain formula klaim logikal \textit{Kanan terminal komputasinya}.
- Terbukti Mutlak!
\end{contoh}

\begin{aktivitas}
  \item Diskusikan dalam studi arsitektur komputasinya: Mengapa pembuktian struktur parameter \textit{Induksi} sangat lekat dalam memvalidasi kebenaran parameter algoritma rekursif (\textit{Recursive Functions/Loops \texttt{while()/for()}})? Paparkan logika abstraksional korelasinya!.
\end{aktivitas}

\begin{asesmen}
Tugas validasi formula iteratif program menggunakan induksi wajib masuk sesi logis evaluasi unjuk kerjanya di Lembar Soal Proyek di depan komisi *Peer Reviewer Student* untuk mendongkrak level keterampilan argumen abstrak (Sub-CPMK 4.2).
\end{asesmen}

\begin{checklist}
  \item Disiplin mematri pemahaman langkah Basis tak boleh salah dan dilewati fungsional cek sirkuit matematisnya. 
  \item Mahir membongkar dan merakit bagian asumsi Hipotesis $P(k)$ logikal untuk membedah dan menata modifikasi fungsional var algoritmik formula rasional target $(k+1)$ pada komputasinya.
\end{checklist}

\end{document}
